% =============================================================================
% ANNEX D : SOFTWARE DEVELOPMENT PLAN (FORMALIZED AND REDACTED VERSION)
% =============================================================================

\chapter{Software Development Plan (formalized and redacted version)}
\label{annex:sdp-markdown}

\begin{tcolorbox}[colback=blue!5, colframe=blue!40!black, boxrule=0.5pt, arc=3pt,
  left=8pt, right=8pt, top=6pt, bottom=6pt]
  \small\itshape
  \textbf{Avertissement.} Le contenu ci-dessous est une reproduction formalisee en \LaTeX{} du Software Development Plan original. Cette version a ete volontairement \textbf{tronquee} et caviardee afin de proteger les informations sensibles (identites, identifiants documentaires internes, liens, details d'infrastructure et metadonnees de pilotage). Des notes de bas de page explicatives ont egalement ete ajoutees pour en faciliter la lecture par un public \textbf{novice}.
\end{tcolorbox}

\section*{1. Document control}

\subsection*{1.1 Roles}
\begin{tabularx}{\textwidth}{|p{3.3cm}|p{2.5cm}|X|p{3cm}|}
\hline
\textbf{Role} & \textbf{Name} & \textbf{Reason for signature} & \textbf{Signature \& date} \\
\hline
System Quality Manager & \caviarder[2.8cm] & Document Owner, Author & \caviarder[2.8cm] \\
\hline
Software Engineer & \caviarder[2.8cm] & Reviewer & \caviarder[2.8cm] \\
\hline
Software Manager & \caviarder[2.8cm] & Approver & \caviarder[2.8cm] \\
\hline
\end{tabularx}

\subsection*{1.2 Version history}
\begin{tabularx}{\textwidth}{|p{2.5cm}|X|}
\hline
\textbf{Version} & \textbf{Description} \\
\hline
1 & Initial version \\
\hline
\end{tabularx}

\subsection*{1.3 References (redacted)}
\begin{longtable}{|p{1.4cm}|p{12cm}|}
\hline
\textbf{Ref.} & \textbf{Document} \\
\hline
R1 & Internal controlled procedure (\caviarder[5.2cm]) \\
\hline
R2 & Internal controlled procedure (\caviarder[5.2cm]) \\
\hline
R3 & Internal controlled procedure (\caviarder[5.2cm]) \\
\hline
R4 & Internal controlled procedure (\caviarder[5.2cm]) \\
\hline
R5 & IEC~62304:2006+AMD1:2015, Medical device software life cycle processes \\
\hline
R6 & Internal controlled procedure (\caviarder[5.2cm]) \\
\hline
R7 & Internal controlled procedure (\caviarder[5.2cm]) \\
\hline
R8 & Internal controlled procedure (\caviarder[5.2cm]) \\
\hline
R9 & Internal controlled procedure (\caviarder[5.2cm]) \\
\hline
R10 & Internal controlled procedure (\caviarder[5.2cm]) \\
\hline
R11 & FDA: Content of Premarket Submissions for Device Software Functions \\
\hline
R12-R19 & Internal project documents (plans, architecture, risk records, detailed design): \caviarder[3.4cm] \\
\hline
\end{longtable}

\subsection*{1.4 Definitions (extract)}
\begin{longtable}{|p{3cm}|p{10.4cm}|}
\hline
\textbf{Term} & \textbf{Definition} \\
\hline
Software System & Integrated collection of Software Items organized to accomplish a function set (IEC~62304 3.28). \\
\hline
Software Item & Any identifiable part of a computer program (IEC~62304 3.25). \\
\hline
Software Unit & Software Item not subdivided into other items (IEC~62304 3.29). \\
\hline
SOUP / OTS & Third-party software not fully controlled through manufacturer life cycle. \\
\hline
SDP & Software Development Plan. \\
\hline
SRS & Software Requirements Specification: defines what software shall do. \\
\hline
SDDD & Software Detailed Design Document: internal design of software items/units. \\
\hline
Software Anomaly & Condition deviating from expected behavior. \\
\hline
\end{longtable}

\section*{2. Purpose and scope}

The SDP defines software development processes for the medical device in compliance with IEC~62304 and internal controlled procedures (\caviarder[2.8cm]).

\subsection*{2.1 Scope}
The plan covers software components in scope: \caviarder[5.5cm].

The document is reviewed at each Gate Review (G1 to G4) and updated under formal change control.

\section*{3. Software safety classification}

\textbf{Preliminary classification:} \caviarder[2.2cm] (conservative assumption for an active medical device).

\textbf{Segregation strategy:}
\begin{itemize}
  \item Critical control loops handled as \caviarder[1.8cm].
  \item Auxiliary applications potentially classified as \caviarder[2.4cm] if architectural segregation is demonstrated.
\end{itemize}

\section*{4. Roles and responsibilities (extrait)}

\begin{tabularx}{\textwidth}{|p{3cm}|X|p{3cm}|}
\hline
\textbf{Role} & \textbf{Responsibilities} & \textbf{Assigned to} \\
\hline
Software Lead & Executes the SDP, arbitrates architecture, and coordinates verification and change impacts. & \caviarder[2.8cm] \\
\hline
Software Developer & Development, unit testing, static analysis, and merge requests. & \caviarder[2.8cm] \\
\hline
Configuration Manager & SCM, CI/CD pipelines, SOUP management, and release baselines. & \caviarder[2.8cm] \\
\hline
Quality Manager & Traceability review, gate validation, and SOP compliance. & \caviarder[2.8cm] \\
\hline
Risk Manager & RMF follow-up\footnote{\textbf{RMF} stands for \textit{Risk Management File}. It contains the records documenting risk analysis, risk control measures, and their verification evidence.}, and safety-impact review of anomalies and changes. & \caviarder[2.8cm] \\
\hline
\end{tabularx}

\section*{5. Development life cycle model}

The project follows a V-Model integrated into the Stage-Gate process.

\begin{tabularx}{\textwidth}{|p{1.2cm}|p{2.2cm}|p{2.5cm}|X|}
\hline
\textbf{Gate} & \textbf{Phase} & \textbf{IEC~62304} & \textbf{Key software activities} \\
\hline
G0 & Planning & 5.1 & User need definition (\caviarder[2.8cm]) \\
\hline
G1 & Input \& Arch. & 5.1, 5.2, 5.3 & SDP, SRS\footnote{\textbf{SRS} stands for \textit{Software Requirements Specification}. This document formalizes verifiable software requirements.}, \caviarder[3.8cm] \\
\hline
G2 & Detailed design & 5.4, 5.5 & SDDD\footnote{\textbf{SDDD} stands for \textit{Software Detailed Design Document}. It decomposes architecture into implementable and testable components/units.}, \caviarder[4.6cm] \\
\hline
G3 & Verification & 5.6, 5.7 & Integration, system testing, \caviarder[2.8cm] \\
\hline
G4 & Validation & 5.8 & Software release and \caviarder[2.4cm] \\
\hline
\end{tabularx}

\medskip
\textbf{Planned deviation:} SDP produced at G1 (instead of G0), since \caviarder[5.8cm].

\section*{6. Traceability and verification framework}

\textbf{Traceability rules:}
\begin{itemize}
  \item Each software requirement has an upstream link to a system requirement and/or risk control measure.
  \item Each software requirement has downstream links to design, code, and verification activity.
  \item Completeness review is performed before each gate.
\end{itemize}

\textbf{Verification strategy (summary):}
\begin{tabularx}{\textwidth}{|p{1.2cm}|p{3.1cm}|X|}
\hline
\textbf{Gate} & \textbf{Verification activity} & \textbf{Acceptance criteria} \\
\hline
G1 & Design review & Requirements and architecture are baselined; \caviarder[3.6cm] \\
\hline
G2 & Unit verification & Static analysis compliant, \caviarder[3.8cm] \\
\hline
G3 & Integration \& system verification & Requirement coverage and integration/system tests validated, \caviarder[2.8cm] \\
\hline
G4 & Release verification & Release package complete; \caviarder[3.4cm] \\
\hline
\end{tabularx}

\section*{7. Software configuration management (SCM)}

SCM is aligned with internal controlled procedures (\caviarder[2.2cm]) using a controlled branching strategy.

\begin{itemize}
  \item \caviarder[2.6cm]: release baseline, direct commits prohibited.
  \item \caviarder[2.6cm]: continuous integration of evolutions.
  \item \caviarder[2.6cm]: preparation of official versions.
  \item \caviarder[2.6cm]: post-release correction with impact analysis.
  \item \caviarder[3.8cm]: development branches.
\end{itemize}

\section*{8. SOUP, cybersecurity, release and maintenance}

\textbf{SOUP management:}
\begin{itemize}
  \item SOUP inventory maintained and periodically reviewed.
  \item SOUP components verified through test plans.
  \item Vulnerability and end-of-life monitoring performed continuously.
\end{itemize}

\textbf{Release:}
release includes version, hash, residual anomalies, SOUP components, and \caviarder[3.8cm], followed by controlled archiving.

\textbf{Problem resolution and maintenance:}
each anomaly follows a documented cycle: logging, triage, investigation, resolution, regression check, then traceability to \caviarder[3.2cm].

\section*{9. Deliverables matrix (redacted summary)}

\begin{tabularx}{\textwidth}{|X|p{2.4cm}|p{1.5cm}|p{2.1cm}|}
\hline
\textbf{Deliverable} & \textbf{Format} & \textbf{Gate} & \textbf{IEC~62304} \\
\hline
SDP & \caviarder[2.2cm] signed & G1 & 5.1 \\
\hline
SRS / Architecture & \caviarder[2.2cm] signed & G1 & 5.2 / 5.3 \\
\hline
SDDD + unit evidence & \caviarder[3.2cm] + PDF signed & G2 & 5.4 / 5.5 \\
\hline
Integration/System tests + traceability matrix & \caviarder[2.6cm] & G3 & 5.6 / 5.7 \\
\hline
Release notes + final SOUP list + software RMF section & \caviarder[2.6cm] & G4 & 5.8 / clause 7 \\
\hline
\end{tabularx}

\medskip
\noindent\textit{Note: detailed document identifiers, storage links, internal audit trails, and operational references are redacted in this portfolio version.}



